% GENERATED FILE. Edit canonical lessons, then run npm run generate:books.
% canonical-source-hash: fnv1a64:21f9532b7721ccba
% canonical-lessons: AR-C511-jura, AR-C511-kabsula, AR-C511-qatra, AR-C511-jabira, AR-C511-dammada, AR-R511-review-504-511

\chapter{Doses and drops}
\label{ch:ar-doses-and-drops}
\hlchaptermodality{511}

\begin{chapteropening}
\textbf{By the end of this chapter:} \emph{I can talk about a dose, a capsule, a drop and a splint, and say that a nurse bound up a wound.}

Complete the last lesson of chapter 511: I can talk about a dose, a capsule, a drop and a splint, and say that a nurse bound up a wound.
\end{chapteropening}

\section[\texorpdfstring{\ar{جرعة} (jurʿa)}{جرعة (jurʿa)}]{\ar{جرعة} (jurʿa) — a dose}

\label{lesson:AR-C511-jura}



\begin{quote}
Before the new one: say the Arabic for to treat, then the Arabic for to cure.
\end{quote}

\subsection*{You'll want to know: \ar{جرعة}}
\textbf{\ar{جرعة}} — \emph{jurʿa} — \textquotedblleft{}a dose\textquotedblright{}.

\begin{quote}
\ar{شرب} \ar{المريض} \ar{جرعة} \ar{من} \ar{الدواء}.
\end{quote}

\begin{quote}
\emph{shariba al-marīḍ jurʿa min ad-dawāʾ.}
\end{quote}

\begin{quote}
The patient drank a dose of the medicine.
\end{quote}

\begin{cousinweb}
\textbf{\ar{جرعة}} (\emph{jurʿa}) is built on the native Arabic root \emph{j-r-ʿ}, \textquotedblleft{}to gulp, to swallow\textquotedblright{} (Lane). Lane gives jurʿa as a single gulp: a dose is one swallow of medicine.
\end{cousinweb}

\begin{grammarlens}[title={What you've built}]
One more word for describing what you lived and felt.
\end{grammarlens}

\subsection*{Your turn}
\begin{itemize}
\raggedright
  \item \emph{Say it:} \emph{jurʿa}
  \item \emph{Say it:} \emph{jurʿa}, once more
  \item \emph{Say it:} \emph{shariba al-marīḍ jurʿa min ad-dawāʾ.} — The patient drank a dose of the medicine.
  \item \emph{Recall:} say the Arabic for to treat, then the Arabic for to cure, then say \emph{jurʿa} again
\end{itemize}

\begin{tcolorbox}[breakable,colback=teal!4,colframe=teal!35!black,title={Before you move on}]
What does \ar{جرعة} mean? (\textquotedblleft{}A dose\textquotedblright{}.) Say it once more.
\end{tcolorbox}
\section[\texorpdfstring{\ar{كبسولة} (kabsūla)}{كبسولة (kabsūla)}]{\ar{كبسولة} (kabsūla) — a capsule}

\label{lesson:AR-C511-kabsula}



\begin{quote}
Before the new one: say the Arabic for to cure, then the Arabic for a dose.
\end{quote}

\subsection*{You'll want to know: \ar{كبسولة}}
\textbf{\ar{كبسولة}} — \emph{kabsūla} — \textquotedblleft{}a capsule\textquotedblright{}.

\begin{quote}
\ar{يأخذ} \ar{جدي} \ar{كبسولة} \ar{كل} \ar{صباح}.
\end{quote}

\begin{quote}
\emph{yaʾkhudhu jaddī kabsūla kull ṣabāḥ.}
\end{quote}

\begin{quote}
My grandfather takes a capsule every morning.
\end{quote}

\begin{cousinweb}
\textbf{\ar{كبسولة}} (\emph{kabsūla}) is a modern loan, usually traced to French \emph{capsule}, from Latin \emph{capsula}, \textquotedblleft{}a small box\textquotedblright{}.
\end{cousinweb}

\begin{grammarlens}[title={What you've built}]
One more word for describing what you lived and felt.
\end{grammarlens}

\subsection*{Your turn}
\begin{itemize}
\raggedright
  \item \emph{Say it:} \emph{kabsūla}
  \item \emph{Say it:} \emph{kabsūla}, once more
  \item \emph{Say it:} \emph{yaʾkhudhu jaddī kabsūla kull ṣabāḥ.} — My grandfather takes a capsule every morning.
  \item \emph{Recall:} say the Arabic for to cure, then the Arabic for a dose, then say \emph{kabsūla} again
\end{itemize}

\begin{tcolorbox}[breakable,colback=teal!4,colframe=teal!35!black,title={Before you move on}]
What does \ar{كبسولة} mean? (\textquotedblleft{}A capsule\textquotedblright{}.) Say it once more.
\end{tcolorbox}
\section[\texorpdfstring{\ar{قطرة} (qaṭra)}{قطرة (qaṭra)}]{\ar{قطرة} (qaṭra) — a droplet; a drop}

\label{lesson:AR-C511-qatra}



\begin{quote}
Before the new one: say the Arabic for a dose, then the Arabic for a capsule.
\end{quote}

\subsection*{You'll want to know: \ar{قطرة}}
\textbf{\ar{قطرة}} — \emph{qaṭra} — \textquotedblleft{}a droplet; a drop\textquotedblright{}.

\begin{quote}
\ar{وضعت} \ar{الممرضة} \ar{قطرة} \ar{في} \ar{عيني}.
\end{quote}

\begin{quote}
\emph{waḍaʿat al-mumarriḍa qaṭra fī ʿaynī.}
\end{quote}

\begin{quote}
The nurse put a drop in my eye.
\end{quote}

\begin{cousinweb}
\textbf{\ar{قطرة}} (\emph{qaṭra}) is built on the native Arabic root \emph{q-ṭ-r}, \textquotedblleft{}to drip, to fall in drops\textquotedblright{} (Lane). Qiṭār, a train, which you know, is from the same root: a line of camels following one another, like drops.
\end{cousinweb}

\begin{grammarlens}[title={What you've built}]
One more word for describing what you lived and felt.
\end{grammarlens}

\subsection*{Your turn}
\begin{itemize}
\raggedright
  \item \emph{Say it:} \emph{qaṭra}
  \item \emph{Say it:} \emph{qaṭra}, once more
  \item \emph{Say it:} \emph{waḍaʿat al-mumarriḍa qaṭra fī ʿaynī.} — The nurse put a drop in my eye.
  \item \emph{Recall:} say the Arabic for a dose, then the Arabic for a capsule, then say \emph{qaṭra} again
\end{itemize}

\begin{tcolorbox}[breakable,colback=teal!4,colframe=teal!35!black,title={Before you move on}]
What does \ar{قطرة} mean? (\textquotedblleft{}A droplet; a drop\textquotedblright{}.) Say it once more.
\end{tcolorbox}
\section[\texorpdfstring{\ar{جبيرة} (jabīra)}{جبيرة (jabīra)}]{\ar{جبيرة} (jabīra) — a splint, a plaster cast}

\label{lesson:AR-C511-jabira}



\begin{quote}
Before the new one: say the Arabic for a capsule, then the Arabic for a droplet.
\end{quote}

\subsection*{You'll want to know: \ar{جبيرة}}
\textbf{\ar{جبيرة}} — \emph{jabīra} — \textquotedblleft{}a splint, a plaster cast\textquotedblright{}.

\begin{quote}
\ar{وضع} \ar{الطبيب} \ar{ذراعي} \ar{في} \ar{جبيرة}.
\end{quote}

\begin{quote}
\emph{waḍaʿa aṭ-ṭabīb dhirāʿī fī jabīra.}
\end{quote}

\begin{quote}
The doctor put my arm in a cast.
\end{quote}

\begin{cousinweb}
\textbf{\ar{جبيرة}} (\emph{jabīra}) is built on the native Arabic root \emph{j-b-r}, \textquotedblleft{}to set a broken bone\textquotedblright{} (Lane). A jabīra holds the set bone in place.
\end{cousinweb}

\begin{grammarlens}[title={What you've built}]
One more word for describing what you lived and felt.
\end{grammarlens}

\subsection*{Your turn}
\begin{itemize}
\raggedright
  \item \emph{Say it:} \emph{jabīra}
  \item \emph{Say it:} \emph{jabīra}, once more
  \item \emph{Say it:} \emph{waḍaʿa aṭ-ṭabīb dhirāʿī fī jabīra.} — The doctor put my arm in a cast.
  \item \emph{Recall:} say the Arabic for a capsule, then the Arabic for a droplet, then say \emph{jabīra} again
\end{itemize}

\begin{tcolorbox}[breakable,colback=teal!4,colframe=teal!35!black,title={Before you move on}]
What does \ar{جبيرة} mean? (\textquotedblleft{}A splint, a plaster cast\textquotedblright{}.) Say it once more.
\end{tcolorbox}
\section[\texorpdfstring{\ar{ضمد} (ḍammada)}{ضمد (ḍammada)}]{\ar{ضمد} (ḍammada) — to bind up (a wound)}

\label{lesson:AR-C511-dammada}



\begin{quote}
Before the new one: say the Arabic for a droplet, then the Arabic for a splint.
\end{quote}

\subsection*{You'll want to know: \ar{ضمد}}
\textbf{\ar{ضمد}} — \emph{ḍammada} — \textquotedblleft{}to bind up (a wound)\textquotedblright{}.

\begin{quote}
\ar{ضمدت} \ar{الممرضة} \ar{جرح} \ar{اللاعب}.
\end{quote}

\begin{quote}
\emph{ḍammadat al-mumarriḍa jurḥ al-lāʿib.}
\end{quote}

\begin{quote}
The nurse bound up the player's wound.
\end{quote}

Present: \ar{يضمد}, \emph{yuḍammid}.

The noun \ar{ضمادة}, a bandage, shares the root.

\begin{cousinweb}
\textbf{\ar{ضمد}} (\emph{ḍammada}) is built on the native Arabic root \emph{ḍ-m-d}, \textquotedblleft{}to bind, to bandage\textquotedblright{} (Lane). Form II ḍammada is to bandage.
\end{cousinweb}

\begin{grammarlens}[title={What you've built}]
One more word for describing what you lived and felt.
\end{grammarlens}

\subsection*{Your turn}
\begin{itemize}
\raggedright
  \item \emph{Say it:} \emph{ḍammada}
  \item \emph{Say it:} \emph{ḍammada}, once more
  \item \emph{Say it:} \emph{ḍammadat al-mumarriḍa jurḥ al-lāʿib.} — The nurse bound up the player's wound.
  \item \emph{Recall:} say the Arabic for a droplet, then the Arabic for a splint, then say \emph{ḍammada} again
\end{itemize}

\begin{tcolorbox}[breakable,colback=teal!4,colframe=teal!35!black,title={Before you move on}]
What does \ar{ضمد} mean? (\textquotedblleft{}To bind up (a wound)\textquotedblright{}.) Say it once more.
\end{tcolorbox}
\section[(dialogue)]{Words from chapters 504–511 — a review}

\label{lesson:AR-R511-review-504-511}



\begin{quote}
Say the Arabic for the liver and the spleen.
\end{quote}

\subsection*{Your turn}
Say these aloud:

\begin{itemize}
\raggedright
  \item the liver, the spleen, the stomach, the intestines, the pulse
  \item a skull, a rib, a thigh, a hip, a joint
  \item an eyelid, an eyelash, the gums, a chin, a nerve
  \item to fall ill, a chill, nausea, to vomit, diarrhoea
  \item a rash, an itch, to scratch, an inflammation, to swell
  \item an infection, a virus, a germ, a vaccination, a swelling
  \item a diagnosis, a lab test, X-rays, to treat, to cure
  \item a dose, a capsule, a droplet, a splint, to bind up
\end{itemize}

\begin{tcolorbox}[breakable,colback=teal!4,colframe=teal!35!black,title={Before you move on}]
The liver? (\textbf{\ar{كبد}}, \emph{kabid}.) The spleen? (\textbf{\ar{طحال}}, \emph{ṭiḥāl}.) The stomach? (\textbf{\ar{معدة}}, \emph{miʿda}.) The intestines? (\textbf{\ar{أمعاء}}, \emph{amʿāʾ}.) The pulse? (\textbf{\ar{نبض}}, \emph{nabḍ}.) A skull? (\textbf{\ar{جمجمة}}, \emph{jumjuma}.) A rib? (\textbf{\ar{ضلع}}, \emph{ḍilʿ}.) A thigh? (\textbf{\ar{فخذ}}, \emph{fakhidh}.) A hip? (\textbf{\ar{ورك}}, \emph{wirk}.) A joint? (\textbf{\ar{مفصل}}, \emph{mafṣil}.) An eyelid? (\textbf{\ar{جفن}}, \emph{jafn}.) An eyelash? (\textbf{\ar{رمش}}, \emph{rimsh}.) The gums? (\textbf{\ar{لثة}}, \emph{litha}.) A chin? (\textbf{\ar{ذقن}}, \emph{dhaqan}.) A nerve? (\textbf{\ar{عصب}}, \emph{ʿaṣab}.) To fall ill? (\textbf{\ar{مرض}}, \emph{mariḍa}.) A chill? (\textbf{\ar{نزلة} \ar{برد}}, \emph{nazlat bard}.) Nausea? (\textbf{\ar{غثيان}}, \emph{ghathayān}.) To vomit? (\textbf{\ar{تقيأ}}, \emph{taqayyaʾa}.) Diarrhoea? (\textbf{\ar{إسهال}}, \emph{ishāl}.) A rash? (\textbf{\ar{طفح} \ar{جلدي}}, \emph{ṭafḥ jildī}.) An itch? (\textbf{\ar{حكة}}, \emph{ḥikka}.) To scratch? (\textbf{\ar{حك}}, \emph{ḥakka}.) An inflammation? (\textbf{\ar{التهاب}}, \emph{iltihāb}.) To swell? (\textbf{\ar{تورم}}, \emph{tawarrama}.) An infection? (\textbf{\ar{عدوى}}, \emph{ʿadwā}.) A virus? (\textbf{\ar{فيروس}}, \emph{fayrūs}.) A germ? (\textbf{\ar{جرثومة}}, \emph{jurthūma}.) A vaccination? (\textbf{\ar{تطعيم}}, \emph{taṭʿīm}.) A swelling? (\textbf{\ar{ورم}}, \emph{waram}.) A diagnosis? (\textbf{\ar{تشخيص}}, \emph{tashkhīṣ}.) A lab test? (\textbf{\ar{تحليل}}, \emph{taḥlīl}.) X-rays? (\textbf{\ar{أشعة}}, \emph{ashiʿʿa}.) To treat? (\textbf{\ar{عالج}}, \emph{ʿālaja}.) To cure? (\textbf{\ar{داوى}}, \emph{dāwā}.) A dose? (\textbf{\ar{جرعة}}, \emph{jurʿa}.) A capsule? (\textbf{\ar{كبسولة}}, \emph{kabsūla}.) A droplet? (\textbf{\ar{قطرة}}, \emph{qaṭra}.) A splint? (\textbf{\ar{جبيرة}}, \emph{jabīra}.) To bind up? (\textbf{\ar{ضمد}}, \emph{ḍammada}.)
\end{tcolorbox}
